\documentclass[11pt]{article}
\usepackage[margin=2.5cm]{geometry}
\usepackage{amsmath,amssymb}
\usepackage{booktabs}
\usepackage{hyperref}

\title{Exact Coordinates for Optimal Heilbronn Configurations\\[4pt]
       \large $n = 5, \ldots, 9$ on the Unit Square}
\author{Companion to: N.\ Sudermann-Merx, \emph{From Computational Certification to Exact Coordinates:
Heilbronn's Triangle Problem on the Unit Square}, 2026.}
\date{}

\begin{document}
\maketitle

\noindent
This document lists the certified optimal coordinates for
Heilbronn's triangle problem on the unit square~$[0,1]^2$
for $n = 5, \ldots, 9$.
All configurations are globally optimal, certified by solving a
mixed-integer quadratically constrained program (MIQCP) with Gurobi
and then deriving exact symbolic coordinates with SymPy.
The derivation is fully automated by the script
\texttt{exact\_coordinates.py} included in this repository.

Points are labeled $1, \ldots, n$.
We denote the optimal minimum triangle area by~$\Delta_n$.

%% ===================================================================
\section*{$n = 5$: \quad $\Delta_5 = \dfrac{\sqrt{3}}{9} \approx 0.19245$}

Four critical triangles. Unique optimal configuration.

\begin{center}
\begin{tabular}{c c c}
\toprule
Pt & $x$ & $y$ \\
\midrule
$1$ & $0$ & $\frac{1}{3}$ \\[4pt]
$2$ & $\frac{\sqrt{3}}{3}$ & $0$ \\[4pt]
$3$ & $1$ & $1 - \frac{\sqrt{3}}{3}$ \\[4pt]
$4$ & $\frac{2}{3}$ & $1$ \\[4pt]
$5$ & $0$ & $1$ \\
\bottomrule
\end{tabular}
\end{center}

%% ===================================================================
\section*{$n = 6$: \quad $\Delta_6 = \dfrac{1}{8} = 0.125$}

Six critical triangles. One-parameter family of optimal solutions
parameterized by $c \in [0, \tfrac{1}{4}]$.

\begin{center}
\begin{tabular}{c c c}
\toprule
Pt & $x$ & $y$ \\
\midrule
$1$ & $0$ & $c$ \\[4pt]
$2$ & $\frac{1}{2}$ & $0$ \\[4pt]
$3$ & $1$ & $\frac{1}{2} - c$ \\[4pt]
$4$ & $\frac{1}{2}$ & $1$ \\[4pt]
$5$ & $0$ & $c + \frac{1}{2}$ \\[4pt]
$6$ & $1$ & $1-c$ \\
\bottomrule
\end{tabular}
\end{center}

%% ===================================================================
\section*{$n = 7$: \quad $\Delta_7 = -\dfrac{a}{2a-4} \approx 0.08386$}

Eight critical triangles. Coordinates are rational functions of
$a \approx 0.2873$, the unique root in $(0,1)$ of the irreducible cubic
\[
  a^3 + 5\,a^2 - 5\,a + 1 = 0.
\]

\begin{center}
\begin{tabular}{c c c}
\toprule
Pt & $x$ & $y$ \\
\midrule
$1$ & $0$ & $0$ \\[4pt]
$2$ & $\dfrac{1}{a^2-3a+2}$ & $0$ \\[8pt]
$3$ & $1$ & $a$ \\[4pt]
$4$ & $\dfrac{2a-1}{a(a-2)}$ & $1$ \\[8pt]
$5$ & $0$ & $1$ \\[4pt]
$6$ & $\dfrac{a-1}{a-2}$ & $a$ \\[8pt]
$7$ & $\dfrac{2(-a^2+3a-1)}{a^3-4a^2+5a-2}$ & $-\dfrac{2a}{a-1}$ \\[8pt]
\bottomrule
\end{tabular}
\end{center}

%% ===================================================================
\section*{$n = 8$: \quad $\Delta_8 = -\dfrac{1}{36}+\dfrac{\sqrt{13}}{36} \approx 0.07238$}

Twelve critical triangles. Point-symmetric configuration with
unique optimal solution involving $\sqrt{13}$.

\begin{center}
\begin{tabular}{c c c}
\toprule
Pt & $x$ & $y$ \\
\midrule
$1$ & $0$ & $0$ \\[4pt]
$2$ & $\frac{1}{6}+\frac{\sqrt{13}}{6}$ & $0$ \\[4pt]
$3$ & $1$ & $\frac{7}{18}-\frac{\sqrt{13}}{18}$ \\[4pt]
$4$ & $1$ & $1$ \\[4pt]
$5$ & $0$ & $\frac{11}{18}+\frac{\sqrt{13}}{18}$ \\[4pt]
$6$ & $\frac{5}{6}-\frac{\sqrt{13}}{6}$ & $1$ \\[4pt]
$7$ & $\frac{5}{6}-\frac{\sqrt{13}}{6}$ & $\frac{7}{9}-\frac{\sqrt{13}}{9}$ \\[4pt]
$8$ & $\frac{1}{6}+\frac{\sqrt{13}}{6}$ & $\frac{2}{9}+\frac{\sqrt{13}}{9}$ \\
\bottomrule
\end{tabular}
\end{center}

%% ===================================================================
\section*{$n = 9$: \quad $\Delta_9 = -\dfrac{11}{64}+\dfrac{9\sqrt{65}}{320} \approx 0.05488$}

Eleven critical triangles. Configuration with four-fold rotational
symmetry; unique optimal solution involving $\sqrt{65}$.

\begin{center}
\begin{tabular}{c c c}
\toprule
Pt & $x$ & $y$ \\
\midrule
$1$ & $0$ & $1-\frac{\sqrt{65}}{10}$ \\[4pt]
$2$ & $\frac{3}{8}-\frac{\sqrt{65}}{40}$ & $0$ \\[4pt]
$3$ & $1$ & $\frac{9}{16}-\frac{3\sqrt{65}}{80}$ \\[4pt]
$4$ & $\frac{3}{8}-\frac{\sqrt{65}}{40}$ & $1$ \\[4pt]
$5$ & $0$ & $\frac{5}{8}+\frac{\sqrt{65}}{40}$ \\[4pt]
$6$ & $\frac{1}{4}+\frac{\sqrt{65}}{20}$ & $\frac{3}{4}-\frac{\sqrt{65}}{20}$ \\[4pt]
$7$ & $\frac{7}{16}+\frac{3\sqrt{65}}{80}$ & $0$ \\[4pt]
$8$ & $\frac{\sqrt{65}}{10}$ & $1$ \\[4pt]
$9$ & $1$ & $\frac{5}{8}+\frac{\sqrt{65}}{40}$ \\
\bottomrule
\end{tabular}
\end{center}

\bigskip\noindent
\textbf{Source:}
\url{https://github.com/spiralulam/heilbronn}

\end{document}
